In comparison to those related works, this dissertation has the following contributions.

\begin{itemize}
    \item 
    In \redd{PART \romannum{1}}, we propose a programming framework for adaptive data analyses where programs represent analysts that can query generalization-preserving mechanisms mediating the access to some data. 
    Based on this  framework, we give formal definition of the notion of adaptivity under the analyst-mechanism model. 
    % \item 

    Then we design a static program analysis algorithm {\THESYSTEM} combining data flow, control flow and  reachability bound analysis in order to provide tight bounds on the adaptivity of a program.
    We also provide a soundness proof of the program analysis showing that the adaptivity estimated by {\THESYSTEM} bounds the \emph{adaptivity} of the program given by the formal definition. 
    
    We also implement this algorithm and reports the experimental evaluation results on several examples.
    % \item An implementation of {\THESYSTEM} and an experimental evaluation of the bounds this implementation provides on several examples.
    \item In \redd{PART \romannum{2}}, we propose
    a path-sensitive reachability-bound algorithm, which computes sound bound on the evaluation times of each program point path-sensitively. This algorithm combines the \emph{amortized bound analysis} through ranking function estimation and the path refinement approach in our algorithm.
    We introduce two novel quantities in this algorithm, the \emph{path reachability-bound} and \emph{loop reachability-bound} and the corresponding estimating algorithms.
    
    Then we implement this algorithm as a prototype and present the evaluation results over four different benchmarks.
  \end{itemize} 